\chapter{Static Analysis}
\label{ch:analysis}\index{static analysis}\index{AdaLang Analyzer}

Chapter~\ref{ch:testing} asks whether the code does the right thing on the inputs we
tried. Chapter~\ref{ch:spark} asks whether it can go wrong on any input at all. A
static analyser asks a third question, and a cheaper one: \emph{does this code contain
shapes that are usually mistakes?} It never runs the program and it proves nothing. It
reads the source, matches patterns a compiler is not obliged to complain about, and
reports them.

That makes it the weakest of the three layers, and the easiest to add. It also makes it
the layer most likely to waste your afternoon, because a pattern that is a mistake in
one codebase is the house idiom in another. This chapter is as much about reading the
output as about running the tool.

\cnote{The closest C equivalents are \texttt{clang-tidy}, \texttt{cppcheck} and the
commercial MISRA checkers. The experience transfers, including the part where the
first run buries one real finding under a thousand style opinions.}

\section{The tool}\index{AdaLang Analyzer!installing}

\textbf{AdaLang Analyzer} is a static analyser for Ada written by
\textbf{Maurizio Martignano} of \textbf{Spazio~IT}. It is free software under the
GPL-3.0. It is built on \textbf{Libadalang}, AdaCore's Ada parsing and semantic
analysis library. That foundation is what lets it resolve names and types, rather
than merely grep for text.

It ships as an Alire crate, so installing it is one line:

\begin{lstlisting}[style=sh]
alr install adalang_analyzer      # lands in ~/.alire/bin
\end{lstlisting}

That build is not quick the first time. Libadalang is a source crate, so a cold install
compiles it from scratch. The binary it produces is self-contained, though, which
matters for continuous integration: see \emph{Running it in CI} below.

The analyser offers around a hundred and twenty-six checks, a few named presets
(\texttt{--recommended}, \texttt{--spark}, \texttt{--automotive}, \texttt{--do178c}),
and an explicit \texttt{-checks=} list. It can also write a baseline of finding
fingerprints and filter against it, which is the feature that turns it from a
one-off audit into something a build can gate on.

\section{The first run, honestly}\index{AdaLang Analyzer!false positives}

Run over the four libraries in this repository, the recommended preset reported
\textbf{953 findings across 381 files}. Not one of them was a real defect.

That is worth sitting with, because it is the normal result and the documentation of
no tool will tell you so. Each high-severity class turned out to be a systematic false
positive with a single identifiable cause:

\begin{itemize}
\item \textbf{31 use-after-free}, every one of them a \lstinline{Free} followed by
  \lstinline{return} or \lstinline{raise}. The analyser does not treat those as ending
  the path, so it reads the next statement as executing after the free.
\item \textbf{8 uninitialised reads}, every one a variable written through an
  \lstinline{Asm} output constraint. Inline assembly is opaque to it.
\item \textbf{90 uninitialised out-parameters}, all the defensive idiom this codebase
  uses everywhere: set the outputs to a safe default, then overwrite them on the
  success path.
\item \textbf{415 swappable parameters}, almost all of them crypto operands declared
  \lstinline{(A, B : Blk)}. Correct in principle, noise here.
\end{itemize}

The lesson is not that the tool is bad. It is that a finding is a \emph{question}, and
the codebase answers it. Read enough of them and the answer becomes obvious in
aggregate, which is faster than triaging one at a time.

\section{The signal is in the empty classes}\index{AdaLang Analyzer!empty rules}

The useful result of that first run was not in what it reported. It was in what it did
not.

Several checks find the errors that actually bite: a condition that is always true, two
case alternatives whose ranges overlap, a boolean operand repeated, an alternative
already covered by an earlier one. Across all four libraries every one of those
reported \textbf{zero}. That is a real statement about the code, and it costs nothing
to keep asking.

There is a catch: those checks are not in the recommended preset. You have to ask for
them by name. A preset is someone else's judgement about which rules matter, and it is
worth overriding.

\section{Choose checks by measuring, not by name}\index{AdaLang Analyzer!choosing checks}

A rule's name tells you what it is for. It does not tell you how it behaves on your
code, and the gap can be enormous.

\lstinline{Exception_Propagation} is the cautionary example. It sounds like exactly
what a bare-metal codebase wants, and under the automotive preset it reports twelve
findings. Enabled on its own it reports \textbf{1,274} in the HAL alone. It fires on
every call that can transitively raise. In a stack that signals errors \emph{with}
exceptions that is very nearly every call: the ext4 layer raises
\lstinline{No_Space}, \lstinline{Corrupt} and \lstinline{Use_Error} by design. It is a restriction,
not a defect detector.

The same discipline rejected four rules that sound like exactly what a readability
check should be:

\begin{itemize}
\item \lstinline{Naming_Convention} reported 1,868 findings, and every single one was
  the same objection: a one-character identifier. That is one house-style disagreement
  repeated, not 1,868 insights.
\item \lstinline{Magic_Number} reported 2,940 in hand-written source across 185 files.
  Register and protocol constants are this codebase's idiom, and the count grows with
  every driver written.
\item \lstinline{No_Multiple_Return} reported 180. Single-exit is a doctrine, not a
  fact; an early return is usually the clearer construct.
\item \lstinline{Long_Line} looked reasonable until the source was measured. The 99th
  percentile is 90 columns and nothing exceeds 120, so every finding came from
  generated SVD files. Formatting belongs to \texttt{gnatformat} anyway.
\end{itemize}

Read the findings before you adopt the rule. Four of the five decisions above reversed
once the output was actually inspected.

\section{What it did find}\index{AdaLang Analyzer!duplicate glue}

One finding paid for the whole exercise. \lstinline{Duplicate_Subprogram} reported that
several peripheral drivers contained the same subprogram body.

They did. Ten drivers had each grown a private copy of the same GPIO-matrix routing
helper: configure the pad as a strong output, then point the matrix output at it. The
analyser found eight of them and missed two more only because those had been named
\lstinline{Route_Out} instead of \lstinline{Drive_Out}. Two of the ten had also
drifted, and did not set \lstinline{OEN_SEL}, so the peripheral's control of its own
output enable depended on a reset default.

The copies are now one \lstinline{ESP32S3.GPIO.Route_Out}
(Chapter~\ref{ch:driver-design}), which also moved the matrix write inside the same
protected object as \lstinline{Configure}. That closed a window where a concurrent
\lstinline{Configure} of the same pad could leave it routed to plain GPIO.

It found a second one, in a much less comfortable place. \lstinline{P256} and
\lstinline{P384} -- the two NIST prime curves, used to authenticate certificate
chains -- had twenty-four subprograms in common, and the bodies were identical
apart from two integer literals. That is exactly what a generic is for, but it is
also SPARK-proven cryptography, so the refactor had to be paid for rather than
asserted: \lstinline{P256} carried 144 proof obligations at zero unproved, and the
only acceptable outcome was the same 144 afterwards.

The arithmetic became two generics, \lstinline{ECC_Bignum} (modular big-integer
arithmetic over a limb count) and \lstinline{ECC_Curve} (Jacobian point arithmetic
for an $a = -3$ short-Weierstrass curve), instantiated at eight limbs and twelve.
The two literals turn out to be the interesting part: \lstinline{Mont_Pow} scanned
the exponent from bit 255 or 383, and \lstinline{Compute_R2} doubled 512 or 768
times. Neither is checkable by the compiler, both are now derived from the limb
count, and getting either wrong in a hand-written copy would have produced a
verifier that silently accepted nothing.

The proof survived unchanged at 144. The unexpected dividend was on the other side:
\lstinline{P384} had no \lstinline{SPARK_Mode} at all, so the curve that verifies
public DoT/DoH roots had been running on the assumption that it matched its proven
sibling. Sharing the source made that assumption testable, and \lstinline{P384}
now proves in its own right -- 135 obligations, zero unproved. GNATprove verifies
generic \emph{instances}, not generics, so this is a genuinely separate result and
needs its own project; what the refactor bought was that the bodies under it are
the ones already known to be right.

The tool is worth having for the eleventh copy, not for the ten it found once.

\section{Baselines, and the two gates}\index{AdaLang Analyzer!baseline}\index{x!analyze}

A backlog of 953 known-benign findings is useless as a gate: nobody reads a report that
is always red. A \emph{baseline} fixes that. The analyser writes a fingerprint for every
current finding, and a later run filters them out and reports only what is new.

The repository drives it through the same command surface as everything else:

\begin{lstlisting}[style=sh]
./x analyze                    # defects
./x analyze --style            # readability, complexity, duplication
./x analyze --update-baseline  # accept the current findings
\end{lstlisting}

The two gates keep separate baselines under \texttt{tools/analyzer-baselines/},
because they enable different checks and neither set of fingerprints can stand in for
the other. Each exits non-zero when a new finding appears.

The fingerprints are not line numbers. Inserting code above a finding does not
invalidate it, which is what makes the baseline survive ordinary editing.

Two things are deliberately \emph{not} baselined. The DO-178C profile is a report
rather than a gate, for reasons in the next section. And the analyser itself is not
part of \texttt{./x test}: it is not in the toolchain this project pins, so a plain
clone would fail a check it cannot run.

\section{Reading the standards profiles}\index{AdaLang Analyzer!profiles}\index{DO-178C}\index{MISRA}

The \texttt{--automotive} and \texttt{--do178c} presets encode industry coding
standards. Both are informative here, and neither is usable as-is.

\texttt{--automotive} reports \textbf{39,511 findings} on this codebase, forty-one
times the recommended set. The bulk is MISRA and AUTOSAR \emph{restrictions}, and a
bare-metal register HAL is built out of precisely what they forbid. There are 25,531
\lstinline{Magic_Number} findings on register offsets alone, plus 1,671 representation
clauses and 136 address clauses. The rest prohibits dynamic allocation, tasking,
controlled types and dispatching, all of which this SDK uses on purpose. Ten of its checks are genuinely about
correctness rather than conformance, and those ten now ride in the default gate. The
rest is a measurement of how far this code sits from a standard it was never written
to.

\texttt{--do178c} is more interesting, and it is honest about itself. Level~A reports
9,551 findings, of which 8,756 are per-subprogram \emph{process} obligations: 1,866
subprograms lack a low-level requirement trace, 5,474 lack an explicit initialiser.
Those grow with every subprogram written, so a baseline over them would need rewriting
on every commit and would signal nothing. Only 120 findings were additive as code
quality.

What the profile does produce is a per-objective \emph{evidence report}, and
\texttt{./x analyze --do178c} writes one per library. It earns its place by stating its
own limits. It says up front that it is verification support, not a compliance
determination. It flags its objective labels as the tool's own paraphrase rather than
DO-178C's normative text. And it lists what it cannot speak to at all: structural
coverage, requirements-based testing, object-code verification, and DO-330 tool
qualification. Treat it as evidence \emph{input} for a certification process, not
as a result.

Two smaller notes on that profile, both non-obvious. Levels~A and~B enable an identical
rule set and differ only in the recorded structural-coverage objective; level~D enables
nothing at all. And two of its rules contradict each other:
\lstinline{No_Compiler_Extensions} flags \lstinline{pragma Loop_Invariant} as
implementation-defined, while \lstinline{Missing_Loop_Variant} demands the matching
\lstinline{pragma Loop_Variant}.

\section{The trap that hides two thirds of your code}\index{AdaLang Analyzer!project profiles}

The analyser loads a GNAT project file, so it sees exactly the sources that project
selects. This repository's HAL scopes \lstinline{Source_Dirs} by runtime profile: under
\texttt{light-tasking} it compiles \texttt{src} and \texttt{svd} only, and under
\texttt{embedded} it takes \texttt{src/**}.

Analysing the default profile therefore covers \textbf{102 of 329 files} and reports a
clean-looking result. The command always passes the widest profile a library supports,
for that reason. If you drive the analyser by hand against a project with scenario
variables, check the \texttt{Files scanned} line before believing the verdict.

\section{Running it in CI}\index{AdaLang Analyzer!continuous integration}

A gate nobody runs is not a gate. The analysis runs as its own job on every pull
request, one step per rule set so the log says which view regressed.

The job caches the \emph{binary}, not the Alire build tree. AdaLang Analyzer links
Libadalang statically, so the installed executable stands alone. It is around
137\,MB, its only dynamic dependency outside the C library is the native GNAT's
\texttt{libgnarl}, and it runs correctly from a copy with the rest of
\texttt{\textasciitilde/.alire} off the path.
Caching one file instead of 1.8\,GB of build output turns a cold run of roughly half an
hour into a restore of about ninety seconds.

\section{Where it sits}\index{verification!layers}

Static analysis is the cheapest layer and the least conclusive. It cannot tell you the
code is right; a clean run means only that the patterns it knows did not match. Its
value is narrow and real: it notices the copy that drifted, the parameter mode that
lies, the identifier that shadows another, and it notices them the day they appear
rather than the day they fail.

Used that way it complements the layers around it rather than competing with them. The
compiler rejects the ill-typed. The tests check the answers. The sweep checks the
foundation. Proof rules out whole classes of failure. And the analyser, sitting
underneath all of them, keeps the shape of the code from quietly drifting while the
other four are busy being right.
